\documentclass[UTF8]{ctexart}
\begin{document}
\section*{AI工具使用详情}
\# AI工具使用详情

本文件按全国大学生数学建模竞赛人工智能工具使用规定撰写。
内容为阶段级摘要，不含 raw prompt、不含题面原文。

- 工具名称：Grok Build 1.0.24
- 控制面：Iteris（jianmo overlay）
- 数值写门：jianmo NumericWriteGate（独立进程复现，非 LLM 盖章）
- 工作区：`B-2026`
- pack\_id：B-2026
- dataset\_digest：sha256:50bf04217eb3e8e7d3ce8b4a84f228b46dcfd93eebae187b443adb7724b2440e
- gate\_pass\_count：6

\#\# 各阶段摘要

\#\#\# S0\_ingest 读题 / 冻结附件
- 是否使用生成式 AI：是
- 核验：人工或程序哈希

\#\#\# S1\_routes 提出正交方法路线
- 是否使用生成式 AI：是
- 核验：人工或程序哈希

\#\#\# S2\_model 机理 / 假设 / 控制方程
- 是否使用生成式 AI：是
- 核验：人工或程序哈希

\#\#\# S3\_numerics 数值实验与写门复现
- 是否使用生成式 AI：是
- 核验：经 NumericWriteGate / 人工核验

\#\#\# S4\_sensitivity 灵敏度
- 是否使用生成式 AI：是
- 核验：经 NumericWriteGate / 人工核验

\#\#\# S5\_paper 论文组装
- 是否使用生成式 AI：是
- 核验：人工或程序哈希

\#\#\# S6\_disclose AI 使用声明
- 是否使用生成式 AI：是
- 核验：人工或程序哈希

\#\# 声明

参赛队对论文全部内容负责。原版本定量数字的来源按上述历史记录保留；本次Q3新增数字来自独立离线测试结果，未经过原版本写门，不应混称同一验证流程。
本文件不包含模型原始对话。

\section*{2026年9月13日 Q3修订补充}
本次使用OpenAI Codex辅助审核Q3证明与代码，发现最近站分界、浮点剪枝、扫描完成状态和内接多边形可行集等问题；辅助提出并推导距离上界半步收缩方法，修改Q3实现与相关论文段落。

验证包括13项回归测试、48个独立离线模型场景及XeLaTeX编译与页面检查。新增结果记录于\texttt{code/q3/offline\_results.json}，修订范围与局限记录于\texttt{Q3\_REVIEW.md}。历史演练数据保留并标注为旧策略，不作为修订代码的官方演练成绩。本次未运行官方模拟器演练或正式测试。

\end{document}
